\section{Discussion}

\input{5.1_physics_interp}

\input{5.2_opt_bc}

\input{5.3_imager_indepent}

\input{5.4_recovery}

\subsection{Limitations}
\label{sec:disc-limitation}


% 5.1 Physical interpretation

% Synthesize, rather than repeat, the land–ocean response, WCO2 sign rule, and
% shadow/brightening bifurcation. Explain why separate land and ocean models are
% physically justified. Distinguish the empirical mechanism evidence from full
% PPDF moment closure.



\input{5.6_future}

% 5.1 Physical interpretation

% Synthesize, rather than repeat, the land–ocean response, WCO2 sign rule, and
% shadow/brightening bifurcation. Explain why separate land and ocean models are
% physically justified. Distinguish the empirical mechanism evidence from full
% PPDF moment closure.

% 5.2 Relationship to the operational bias correction

% Discuss the raw/BC/ML experiment:

% - an ML model trained from raw retrievals recovers much of the operational
%   increment;
% - ML applied to bias-corrected XCO2 performs best overall;
% - near-cloud land is complementary: the operational correction can
%   over-correct, while the learned residual model partly reverses it.

% Frame the proposed method as a complementary residual correction, not a
% replacement for the operational correction.

% 5.3 Meaning and limits of imager independence

% Use three tiers:

% 1. **Deployment:** no imager or neighboring footprints at inference.
% 2. **Sensitivity:** the spectrum contains cloud-proximity information,
%    including a response below the MODIS resolution floor.
% 3. **Transferability:** the conceptual requirements are resolved absorption
%    bands, channel-level prior optical depth, and single-footprint spectra.

% State the caveat once:

% > The imager serves as scaffolding for target construction and validation but
% > is removed from the deployed correction.

% Discuss post-2022 OCO-2 use and OCO-3, CO2M, GOSAT-GW, and TEMPO as prospects,
% not demonstrated cross-sensor equivalence.

% 5.4 Observation-recovery implications

% Use production QF1 counts to quantify candidate recovery by distance, surface,
% and region. Avoid calling observations “usable” solely because RMSE is lower;
% formal acceptance also requires retrieval validity, uncertainty calibration,
% and downstream application tests.

% Recommended language is **candidates for recovery** until an explicit
% acceptance criterion and inversion experiment exist.

% 5.5 Limitations

% Collect the limitations in one subsection:

% - anomaly-target circularity and possible real-gradient contamination;
% - MYD35 mask semantics, sub-pixel cloud, and parallax limitations;
% - land-heavy TCCON sampling and limited ocean-reference sample;
% - high-latitude and representation-error residuals;
% - absence of fitted-versus-tallied PPDF moment closure;
% - possible attenuation of local CO2 contrast by retrieval-state predictors.

% #### 5.6 Future work

% Keep this short: Monte Carlo PPDF closure, plume-injection OSSE,
% transport-model decomposition of the target, cross-sensor validation, and a
% downstream flux-inversion assessment.